\scanpage{97}

Sous ces conditions, $B$ est étale sur $A$.

On peut supposer [quitte à faire une
\deleted{\unsure{extens.}} changmt de base $A'\unsure{-|\to} A$, $A'$ la
clôture \unsure{strictemt} locale de $A$] que $A$ est
\unsure{strictemt} local. \deleted{Alors il est \unsure{connu}} De plus
on peut supposer $B$ intègre. Alors il
est connu \inserted{(Abhyankar)}, grâce à (i) \inserted{(ii)}, que
$\Gamma=\mathbb Z/n\mathbb Z$, $B$ est de la forme $A[T]/(T^n-x)$.

% slanted
\margin{\uline{expliciter en} \uline{Lemme}}

Si $A'=A/yA$, alors $B\otimes_A A'=A'[T]/(T^n-x')$,
où $x'$ est l'image de $x$ dans $A'$, qui est un
uniformisant. On conclut \unsure{maintent|maintenant}
grâce à (iii).

% slanted
\margin{\deleted{N.B. On a laissé tomber dan\unsure{s}}}

\deleted{\unsure{Coroll.|Corollaire}} \uline{Théorème} \quad Soit \deleted{$A$ un anneau noeth.}
\deleted{$f:X\longrightarrow Y$ un morphisme essentiellement}
$f:X\longrightarrow Y$ un morphisme lisse,
\deleted{$Y$ loc. \unsure{noeth}}. Alors $f$ est loc.
1-asphérique pour les groupes finis
d'ordre premier aux car. résiduelles
de $Y$.

% archive sheet number: 96
